% !PRINTED 305
\weberpage{305}

The problem of finding all these metacyclic equations is
thus reduced to finding all metacyclic divisors of the homogeneous
congruence group $S$.

If we take as an example the case $p = 3$, $k = 2$, and thus
ask for the metacyclic equations of degree 9, then by this
theorem we are concerned with the group $S$ of the linear
substitutions taken modulo 3

\weq{12}{\sigma = \begin{pmatrix} a, & b \\ c, & d \end{pmatrix},}

with which we occupied ourselves, from a somewhat different
point of view, in \S.\,66.

There we first considered a divisor $E$ of $S$, which consists
of all substitutions $\sigma$ with the condition $a\,d - b\,c$
$\equiv 1$ (mod 3), and moreover combined two substitutions

\weq{13}{\begin{pmatrix} a, & b \\ c, & d \end{pmatrix}, \quad \begin{pmatrix} -a, & -b \\ -c, & -d \end{pmatrix}}

into a single element. The group so specialised was of
degree 12 and had a normal divisor of index 3:

\[ G = \begin{pmatrix} 1, & 0 \\ 0, & 1 \end{pmatrix}, \begin{pmatrix} 0, & 1 \\ -1, & 0 \end{pmatrix}, \begin{pmatrix} -1, & 1 \\ 1, & 1 \end{pmatrix}, \begin{pmatrix} 1, & 1 \\ 1, & -1 \end{pmatrix}, \]

from which $E$ was composed thus:

\[ E = G + \begin{pmatrix} 1, & 0 \\ 1, & 1 \end{pmatrix} G + \begin{pmatrix} 1, & 0 \\ -1, & 1 \end{pmatrix} G. \]

In the group $S$, however, the two substitutions (13)
count as distinct, and moreover the substitutions
$\sigma$ whose determinant is $a\,d - b\,c \equiv -1$ (mod 3) must be
added, whereby the degree of the group rises to 48.
$G$ is extended, by the change of the signs of
all its elements, to a group of degree 8, and $E$ to a
group of degree 24, of which the extended $G$ is normal divisor.
Here the whole group $S$ is now \gesperrt{metacyclic}, as one
sees from the following composition, in which the extended
group $G$ is denoted by $S_2$ and the extended group $E$ by $S_1$:

\[ S_4 = \begin{pmatrix} 1, & 0 \\ 0, & 1 \end{pmatrix}, \begin{pmatrix} -1, & 0 \\ 0, & -1 \end{pmatrix} \]
%CONT>>
